\section{Discussion}
\label{sec:discussion}

\textbf{Limitations.}
(1)~Evidence. The decision results come from simulation, model-based or with measured parameters. The generative model departs from the planner's assumptions on purpose,
but it is ours. \RealDataLimitation{} Labels in the measured runs are simulated draws from laws estimated from 8--101 trials, so a tail beyond what we saw (up to 2.9\% on
Element) cannot be excluded, and it would either break the budget or make \sys ask.
(2)~Cross-family dependence. Product certification assumes A2. A stale history violates it in the model (\S\ref{sec:eval-pop}). On Element no two families were wrong
together in \RealCoN{} trials, which bounds co-failure only at a few percent. Joint certification removes the assumption but needs about $\ln(1/\delta)/\alpha_{\mathrm{test}}$ labels
per stratum and fault class, which only pooling provides.
(3)~Calibration population. Bounds hold only for labels from interrupted effects. Passive labels, an unlogged lag tail, or an undetected delayed recovery in Doze break the
budget (Tables~\ref{tab:realsim},~\ref{tab:stress}). Drift within a version is not detected.
(4)~Escalations. Joint certification escalates two thirds of interrupted workflows at 2{,}000 labels in the model and every interrupted K-9 send. Each ask interrupts a
user who delegated the task. We have not run a user study of the prompt (it needs IRB approval).
(5)~Effects. A missed commit point is invisible to \sys, and an app that commits half a batch needs app-specific probes.

\textbf{Android realism.}
Payment apps resist this agent. \texttt{FLAG\_SECURE} blocks screenshots. Since Android~14, \texttt{accessibility\-Data\-Sensitive} views are hidden from services that do
not declare \texttt{isAccessibilityTool}~\cite{android_a11ysensitive}. Google Play reserves that declaration for assistive tools~\cite{googleplay_a11ypolicy}. Banking apps
may also detect accessibility services and gate the commit behind BiometricPrompt. Our measurements therefore use open-source apps with self-hosted servers. Payments need
MockBank, a test app we plan to build. Element, Simple SMS and Thunderbird offer only counterparty-visible nonce fields (Table~\ref{tab:apps}), and our Element runs opted in.
Physical phones kill and defer more aggressively, so lag tails may exceed 13.6\,s. A withheld reply took another Element code path and stayed exactly once in
\RealTimeoutClean{} clean trials. End-to-end encrypted rooms are untested.

\textbf{Ethics and threat model.}
A nonce writes agent text into user content, so \sys uses user-only fields by default, requires opt-in for counterparty-visible ones, and never uses invisible characters. We assume
a benign agent and apps and an untrusted co-installed app that may race for the foreground or render look-alikes. Foreground re-verification narrows, but does not close, the
rebinding window~\cite{qian2026rebinding}, and only outgoing items of the verified account count as nonce matches. The model applies to desktop and web agents~\cite{xie2024osworld,zhou2024webarena}; the two-hop reduction is where phones differ.

\textbf{Data and artifact availability.}
Every table, figure and inline number is generated from the artifact's CSVs by a standard-library Python generator: the seeded simulator, calibration, ledger benchmark, and the
emulator measurements with per-trial CSVs and logs, plus the rig (AVD configuration, APK hashes, fault proxy, driver, reader). An anonymized copy will be provided to reviewers.
